\section{Conclusion}
\label{sec:conclusion}

The discounted multi-step objective extends goal-frontier maximization from
myopic action evaluation to temporal planning without sacrificing the
self-balancing property. Theorem~\ref{thm:discounted_self_balancing}
establishes that the structural guarantees of the foundational paper hold at
every time step and in the aggregate: sustained destruction, coercion, and
rigidity are value-negative under any discount factor.

This extension enables three capabilities the myopic framework cannot
support. First, risk-capability identification: the actor can evaluate
dangerous capability combinations through causal exercise pathways and
justify preemptive restriction when the discounted expected contraction
exceeds the immediate restriction cost
(Proposition~\ref{prop:preemptive}). Second, a criterion under which the structural-avoidance incentive is counteracted:
the actor accepts short-term $\volL$-contraction from structural discovery
because the information value of accurate structure improves future planning
quality (Proposition~\ref{prop:structure_value}). Third, the anti-monopolar
property: under a diversity condition (positive external discovery rate),
full capability domination is anti-maximizing because it collapses the
external growth that drives future $\volL$ expansion
(Proposition~\ref{prop:anti_monopolar}). Two formal mechanisms support the
exclusion. Observational individuation
(Corollary~\ref{cor:a1_from_oi}) restores static protection: every
distinguishable agent contributes at least one irreplaceable poset node,
making each elimination step $\volL$-contracting even under full individual
subsumption. The gradual domination result
(Corollary~\ref{cor:gradual}) provides dynamic protection: when domination
takes more than $\tau_D^*$ steps and non-members contribute ongoing novel
capabilities, the diversity trajectory dominates at any discount factor.
Together, these address the foundational paper's substitution problem: a
GFM-aligned agent in a population of distinguishable agents faces both a
per-step cost and a trajectory-level disadvantage from pursuing domination.
The result creates endogenous anti-singleton pressure under conditions that
hold generically, though it does not constitute unconditional exclusion:
the dynamic protection requires ongoing external novelty and a
linearized-growth approximation whose accuracy is ultimately empirical.

The verification asymmetry in risk communication remains a partial resolution.
Causal-counterfactual evaluation narrows the gap by basing risk-trust on the
structural plausibility of claims rather than on observed outcomes, but it
cannot fully close it: a plausible risk whose probability is genuinely
uncertain remains in trust limbo until either the catastrophe occurs or
enough evidence accumulates to revise the probability estimate. The Wamura
problem---that successful prevention destroys the evidence that would have
justified it---is an inherent feature of risk reasoning, not a deficiency of
this framework.

Together with the foundational paper's structural alignment properties and the
companion paper's computable poset and leverage dynamics, the discounted
objective completes the theoretical ingredients for a GFM actor: a
computable measure, an endogenous importance signal, and a temporal objective
that supports planning, risk identification, and structural investment. Open
questions remain at both the theoretical level (subsumption correction,
intrinsic value, optimal $\alpha$) and the engineering level (planning
algorithms, pathway enumeration, discount-factor selection). The framework
provides the formal foundation on which implementation can proceed; the
remaining gaps are in approximate-planning methods and empirical validation
in a concrete environment, which are the subjects of planned follow-up work.
